% -*- coding: utf-8 -*-
% !TEX program = xelatex
\documentclass[plain,noanswer,medmath]{../../template/jnuexam}
\usepackage{../../template/kysx}

\begin{document}


\examtitle{2011}{3}


\loadproblems{1}{2011P1}%载入试卷一
\loadproblems{2}{2011P2}%载入试卷二


\makepart{选择题}{1～8小题，每小题4分，共32分}


\useproblem{2}{1}{1}%【同试卷二第一(1)题】


\useproblem{2}{1}{2}%【同试卷二第一(2)题】


\begin{problem}
设$\left\{u_n \right\}$是数列，则下列命题正确的是\pickout{}
\begin{abcd}
\item 若$\sum_{n = 1}^{\infty}u_n$收敛，则$\sum_{n = 1}^{\infty}(u_{2n - 1} + u_{2n})$收敛．
\item 若$\sum_{n = 1}^{\infty}(u_{2n - 1} + u_{2n})$收敛，则$\sum_{n = 1}^{\infty}u_n$收敛．
\item 若$\sum_{n = 1}^{\infty}u_n$收敛，则$\sum_{n = 1}^{\infty}(u_{2n - 1} - u_{2n})$收敛．
\item 若$\sum_{n = 1}^{\infty}(u_{2n - 1} - u_{2n})$收敛，则$\sum_{n = 1}^{\infty}u_n$收敛．
\end{abcd}
\end{problem}


\useproblem{1}{1}{4}%【同试卷一第一(4)题】


\useproblem{1}{1}{5}%【同试卷一第一(5)题】


\begin{problem}
设$A$为$4 \times 3$矩阵，$\eta_1,\eta_2,\eta_3$是非齐次线性方程组$Ax = \beta$的$3$个线性无关的解，
$k_1,k_2$为任意常数，则$Ax = \beta$的通解为\pickout{}
\begin{abcd}
\item $\frac{\eta_2 + \eta_3}{2} + k_1(\eta_2 - \eta_1)$．
\item $\frac{\eta_2 - \eta_3}{2} + k_1(\eta_2 - \eta_1)$．
\item $\frac{\eta_2 + \eta_3}{2} + k_1(\eta_2 - \eta_1) + k_2(\eta_3 - \eta_1)$．
\item $\frac{\eta_2 - \eta_3}{2} + k_1(\eta_2 - \eta_1) + k_2(\eta_3 - \eta_1)$．
\end{abcd}
\end{problem}


\useproblem{1}{1}{7}%【同试卷一第一(7)题】


\begin{problem}
设总体$X$服从参数为$\lambda$（$\lambda > 0$）的泊松分布，
$X_1,X_2, \cdots ,X_n$（$n \ge 2$）为来自总体$X$的简单随机样本，
则对于统计量$T_1 = \frac{1}{n}\sum_{i = 1}^n X_i$，
$T_2 = \frac{1}{n - 1}\sum_{i = 1}^{n - 1} X_i + \frac{1}{n} X_n$，有\pickout{}
\begin{abcd}
\item $E(T_1) > E(T_2)$，$D(T_1) > D(T_2)$．
\item $E(T_1) > E(T_2)$，$D(T_1) > D(T_2)$．
\item $E(T_1) < E(T_2)$，$D(T_1) < D(T_2)$．
\item $E(T_1) < E(T_2)$，$D(T_1) < D(T_2)$．
\end{abcd}
\end{problem}


\makepart{填空题}{9～14小题，每小题4分，共24分}


\begin{problem}
设$f(x) = \lim_{t \to 0} x(1 + 3t)^{\frac{x}{t}}$，
则$f'(x) =$ \fillin{}．
\end{problem}


\begin{problem}
设函数$z = \left(1 + \frac{x}{y} \right)^{\frac{x}{y}}$，则$\dz\big|_{(1,1)} =$ \fillin{}．
\end{problem}


\begin{problem}
曲线$\tan \left(x + y + \frac{\pi}{4} \right) = \e^{y}$在点$\left(0,0 \right)$处的切线方程为 \fillin{}．
\end{problem}


\begin{problem}
曲线$y = \sqrt{x^{2} - 1}$，直线$x = 2$及$x$轴所围成的平面图形绕$x$轴旋转所成的旋转体的体积为 \fillin{}．
\end{problem}


\begin{problem}
设二次型$f\left(x_{1},\,x_{2},\,x_{3} \right) = x^T Ax$的秩为$1$，$A$的各行元素之和为$3$，
则$f$在正交变换$x = Qy$下的标准形为 \fillin{}．
\end{problem}


\useproblem{1}{2}{14}%【同试卷一第二(14)题】


\makepart{解答题}{15～23小题，共94分}


\begin{problem}[points=10]
求极限$\lim_{x \to 0} \frac{\sqrt{1 + 2\sin x} - x - 1}{x\ln \left(1 + x \right)}$．
\end{problem}


\begin{problem}[points=10]
已知函数$f(u,v)$具有连续的二阶偏导数，$f(1,1) = 2$是$f(u,v)$的极值，
$z = f[x+y, f(x,y)]$，求$\frac{\pd^2 z}{\pdx\pd y}\big|_{(1,1)}$．
\end{problem}


\begin{problem}[points=10]
求$\int \frac{\arcsin \sqrt{x} + \ln x}{\sqrt{x}} \dx$．
\end{problem}


\begin{problem}[points=10]
证明$4\arctan x - x + \frac{4\pi}{3} - \sqrt{3} = 0$恰有两个实根．
\end{problem}


\begin{problem}[points=10]
设函数$f(x)$在$[0,1]$有连续导数，$f(0) = 1$，且
$$\iint_{D_t} f'(x + y)\dx\dy = \iint_{D_t} f(t)\dx\dy,$$ 
$D_t = \left\{(x,y)\left| 0 \le y \le t - x,0 \le x \le t \right. \right\}$（$0 < t \le 1$），
求$f(x)$的表达式．
\end{problem}


\useproblem[points=11]{1}{3}{20}%【同试卷一第三(20)题】


\useproblem[points=11]{1}{3}{21}%【同试卷一第三(21)题】


\useproblem[points=11]{1}{3}{22}%【同试卷一第三(22)题】


\begin{problem}[points=11]
设二维随机变量$(X,Y)$服从区域$G$上的均匀分布，
其中$G$是由$x - y = 0$，\goodbreak$x + y = 2$与$y = 0$所围成的区域．\par
\begin{enumerate*}
\item 求边缘概率密度$f_X(x)$；
\item 求条件密度函数$f_{X|Y}(x|y)$．
\end{enumerate*}
\end{problem}


\end{document}
